\documentclass[11pt,a4paper]{article}
\usepackage[utf8]{inputenc}
\usepackage{lmodern}
\usepackage{amsmath, amssymb, amsthm, graphicx, float}
\usepackage[margin=1in]{geometry}
\graphicspath{{../../}}

% Define formal mathematical environments
\newtheorem{theorem}{Theorem}[section]
\newtheorem{proposition}[theorem]{Proposition}
\newtheorem{lemma}[theorem]{Lemma}
\newtheorem{definition}[theorem]{Definition}
\newtheorem{remark}[theorem]{Remark}

\title{Discrete Almgren-Chriss Optimal Execution}
\author{Sacha Reymond}
\date{September 28, 2026}

\begin{document}

\maketitle

\begin{abstract}
\noindent \textbf{Project Objective:} This paper presents a discrete-time Almgren-Chriss optimal execution engine extended with a dynamic Machine Learning Model Predictive Control (MPC) framework. The system leverages localized order book features to forecast short-term liquidity and volatility, dynamically adapting the liquidation schedule to navigate microstructural risk. \\

\noindent \textbf{Technical Stack:} The computational engine is developed in Python, utilizing NumPy for strict vectorization, SciPy's banded solver (\texttt{scipy.linalg.solve\_banded}) for $O(N)$ tridiagonal matrix inversion, and Scikit-learn for regularized Ridge regression and tree-based predictive modeling. \\

\noindent \textbf{Core Findings:} Theoretical validation confirms that the discrete numerical solver aligns with continuous-time analytical limits to within $0.11\%$ discretization error. Furthermore, empirical historical bootstrapping demonstrates that the dynamic MPC model successfully identifies and adapts to non-linear liquidity vacuums, safely reverting to the theoretically optimal variance-minimizing schedule during periods of severe market stress.
\end{abstract}

\section{Model Setup and Market Dynamics}

\begin{definition}[Trading Schedule]
Let $T$ be the total time horizon and $N \in \mathbb{N}$ the number of discrete trading intervals. We define the time step $\tau = \frac{T}{N}$ and the discrete time grid $t_k = k \tau$ for $k \in \{0, 1, \dots, N\}$. 

Let $X$ denote the initial inventory. The inventory at time $t_k$ is denoted $x_k$, subject to the boundary conditions $x_0 = X$ and $x_N = 0$. The number of shares liquidated between $t_{k-1}$ and $t_k$ is given by:
\[ v_k = x_{k-1} - x_k \quad \text{for} \quad k \in \{1, \dots, N\} \]
\end{definition}

\begin{definition}[Price Dynamics and Execution Cost]
The fundamental price process follows an arithmetic random walk with permanent market impact:
\[ S_k = S_{k-1} + \sigma \sqrt{\tau} \xi_k - \gamma v_k \]
where $\xi_k \overset{iid}{\sim} \mathcal{N}(0, 1)$, $\sigma$ is the arithmetic volatility, and $\gamma \geq 0$ is the permanent impact parameter.

The actual execution price incorporates a temporary market impact penalty $\eta \geq 0$:
\[ \tilde{S}_k = S_{k-1} - \eta \frac{v_k}{\tau} \]
The total cost of execution is the sum of cash inflows:
\[ C(x) = \sum_{k=1}^{N} v_k \tilde{S}_k \]
\end{definition}

\section{Implementation Shortfall and Objective Function}

\begin{definition}[Implementation Shortfall]
The Implementation Shortfall (IS) measures the execution cost relative to the initial paper value of the portfolio:
\[ IS = C(x) - X S_0 \]
\end{definition}

\begin{proposition}[Moments of the Implementation Shortfall]
The expected value and variance of the Implementation Shortfall are given by:
\begin{align*}
E[IS] &= - \sum_{k=1}^{N} \left( v_k \gamma (x_0 - x_{k-1}) + \eta \frac{v_k^2}{\tau} \right) \\
V[IS] &= \sigma^2 \tau \sum_{i=1}^{N-1} x_i^2
\end{align*}
\end{proposition}

\begin{proof}
The Implementation Shortfall is given by:
\[ IS = \sum_{k=1}^{N} \left( v_k S_{k-1} - \eta \frac{v_k^2}{\tau} \right) - X S_0 \]
Substituting the unrolled fundamental price $S_{k-1} = S_0 + \sigma \sqrt{\tau} \sum_{i=1}^{k-1} \xi_i - \gamma(x_0 - x_{k-1})$, the expected value extracts the deterministic components, yielding $E[IS]$.

For the variance, we isolate the stochastic term of $IS$:
\[ IS_{\text{stochastic}} = \sigma \sqrt{\tau} \sum_{k=1}^{N} v_k \sum_{i=1}^{k-1} \xi_i \]
By swapping the order of summation over the discrete indices from $\{(k,i) : 1 \le k \le N, 1 \le i \le k-1\}$ to $\{(i,k) : 1 \le i \le N-1, i+1 \le k \le N\}$, we obtain:
\[ IS_{\text{stochastic}} = \sigma \sqrt{\tau} \sum_{i=1}^{N-1} \xi_i \sum_{k=i+1}^{N} v_k \]
Since the sum of executed volumes from $i+1$ to $N$ is exactly the remaining inventory $x_i - x_N$, and $x_N = 0$, this simplifies to $\sigma \sqrt{\tau} \sum_{i=1}^{N-1} x_i \xi_i$. Because $\xi_i$ are independent standard normal variables, the variance evaluates to the sum of the squared coefficients:
\[ V[IS] = \sigma^2 \tau \sum_{i=1}^{N-1} x_i^2 \]
\end{proof}

\begin{definition}[Mean-Variance Objective Function]
Because $IS$ represents a negative loss, we aim to minimize the expected execution penalty $E = -E[IS]$ penalized by the strategy's variance. For a risk aversion parameter $\lambda \geq 0$, the objective function is:
\[ U(x) = E + \lambda V[IS] \]
\end{definition}

\begin{proposition}[Quadratic Form of the Objective Function]
The objective function can be isolated with respect to the interior points $x = (x_1, x_2, \dots, x_{N-1})$ as a strict multi-dimensional quadratic form $\tilde{U}(x) = \frac{1}{2} x^T A x - b^T x$, where $\tilde{U}(x)$ is defined as:
\[ \tilde{U}(x) = \frac{\tilde{\eta}}{\tau} \sum_{k=1}^{N} (x_{k-1} - x_k)^2 + \lambda \sigma^2 \tau \sum_{k=1}^{N-1} x_k^2 \qquad \text{with } \tilde{\eta} = \eta - \frac{1}{2} \gamma \tau \]
The elements of the symmetric tridiagonal matrix $A$ and vector $b$ are:
\begin{align*}
A_{k,k} &= \frac{4\tilde{\eta}}{\tau} + 2 \lambda \sigma^2 \tau \\
A_{k, k+1} &= A_{k+1, k} = -\frac{2\tilde{\eta}}{\tau} \\
b_1 &= \frac{2\tilde{\eta}X}{\tau}, \quad b_k = 0 \quad \text{for } k > 1
\end{align*}
\end{proposition}

\begin{proof}
We expand the impact penalty sum $\sum_{k=1}^{N} (x_{k-1} - x_k)^2$:
\[ \sum_{k=1}^{N} (x_{k-1} - x_k)^2 = x_0^2 - 2x_0 x_1 + 2\sum_{k=1}^{N-1} x_k^2 - 2\sum_{k=1}^{N-2} x_k x_{k+1} + x_N^2 \]
Since $x_0 = X$ and $x_N = 0$, substituting this into $\tilde{U}(x)$ yields:
\[ \tilde{U}(x) = \frac{\tilde{\eta}}{\tau} X^2 - \frac{2\tilde{\eta}}{\tau} X x_1 - \frac{2\tilde{\eta}}{\tau} \sum_{k=1}^{N-2} x_k x_{k+1} + \sum_{k=1}^{N-1} \left( \frac{2\tilde{\eta}}{\tau} + \lambda \sigma^2 \tau \right) x_k^2 \]
By mapping this to the standard quadratic form $\frac{1}{2}x^T A x - b^T x$, we multiply the $x_k^2$ coefficients by $2$ to identify the diagonal elements $A_{k,k} = \frac{4\tilde{\eta}}{\tau} + 2\lambda \sigma^2 \tau$, identify the off-diagonal elements $A_{k,k+1} = -\frac{2\tilde{\eta}}{\tau}$, and extract the linear vector component $b_1 = \frac{2\tilde{\eta}X}{\tau}$.
\end{proof}

\begin{theorem}[Strict Convexity and Global Minimum]
The matrix $A = \nabla^2 \tilde{U}(x)$ is positive definite, meaning for all non-zero vectors $x \in \mathbb{R}^{N-1}$, $x^T A x > 0$. Consequently, $\tilde{U}(x)$ is strictly convex and admits a unique global minimum $x^*$.
\end{theorem}

\begin{proof}
For any non-zero vector $x \in \mathbb{R}^{N-1}$, the quadratic form expands as:
\[ x^T A x = \left(\frac{4\tilde{\eta}}{\tau} + 2\lambda\sigma^2\tau\right)\sum_{k=1}^{N-1} x_k^2 - \frac{4\tilde{\eta}}{\tau} \sum_{k=1}^{N-2} x_k x_{k+1} \]
Factoring out $\frac{2\tilde{\eta}}{\tau}$, we rewrite the sum:
\[ x^T A x = \frac{2\tilde{\eta}}{\tau} \left( 2\sum_{k=1}^{N-1} x_k^2 - 2\sum_{k=1}^{N-2} x_k x_{k+1} \right) + 2\lambda\sigma^2\tau \sum_{k=1}^{N-1} x_k^2 \]
Applying the algebraic identity $(a-b)^2 = a^2 - 2ab + b^2$ allows us to regroup the bracketed terms:
\[ x^T A x = \frac{2\tilde{\eta}}{\tau} \left[ x_1^2 + x_{N-1}^2 + \sum_{k=1}^{N-2} (x_k - x_{k+1})^2 \right] + 2\lambda\sigma^2\tau \sum_{k=1}^{N-1} x_k^2 \]
Since $\tilde{\eta}, \tau, \lambda, \sigma^2 > 0$, every term is a non-negative sum of squares. For $x \neq 0$, at least one coordinate is non-zero, ensuring $\sum x_k^2 > 0$. Thus, $x^T A x > 0$.
\end{proof}

\begin{proposition}[Risk-Neutral Execution Limit (TWAP)]
As the risk aversion parameter $\lambda \to 0$, the optimal execution strategy converges to the Time-Weighted Average Price (TWAP) schedule, where the executed volume at each step is strictly constant: $v_k = \frac{X}{N}$ for all $k \in \{1, \dots, N\}$.
\end{proposition}

\begin{proof}
For a risk-neutral agent ($\lambda = 0$), the objective function reduces entirely to the temporary market impact penalty:
\[ \min_{\mathbf{v}} \sum_{k=1}^N v_k^2 \]
subject to the liquidation constraint $\sum_{k=1}^N v_k = X$.

Let $f(v) = v^2$. Since the second derivative $f''(v) = 2 > 0$, the function is strictly convex. Applying Jensen's inequality with uniform weights $w_k = \frac{1}{N}$:
\[ f\left( \frac{1}{N} \sum_{k=1}^N v_k \right) \le \frac{1}{N} \sum_{k=1}^N f(v_k) \]
Substituting the constraint $\sum_{k=1}^N v_k = X$ into the left side yields:
\[ \left( \frac{X}{N} \right)^2 \le \frac{1}{N} \sum_{k=1}^N v_k^2 \implies \sum_{k=1}^N v_k^2 \ge \frac{X^2}{N} \]
Because $f(v)$ is strictly convex, Jensen's inequality holds as a strict equality if and only if all variables are identical. Therefore, the absolute minimum is uniquely achieved when $v_1 = v_2 = \dots = v_N = \frac{X}{N}$.
\end{proof}

\begin{remark}[High Risk Aversion Boundary]
Conversely, as $\lambda$ increases, the variance penalty $\sum_{k=1}^{N-1} x_k^2$ dominates the objective function. To minimize this holding area, the optimization algorithm generates a highly convex, front-loaded schedule, aggressively liquidating inventory in the earliest trading intervals to shed stochastic risk.
\end{remark}

\begin{figure}[htbp]
    \centering
    \includegraphics[width=0.7\textwidth]{results/Optimal_Execution_Trajectories_across_Risk_Aversion_Levels.png}
    \caption{We implemented the Almgren-Chriss model for 50 different values of $\lambda$ ranging from $10^{-6}$ to $10$. We can clearly observe how the curvature evolves as $\lambda$ gets bigger and how it tends to become a straight line as $\lambda \to 0$.}
    \label{fig:mon_image} 
\end{figure}

\section{Continuous-Time Limit and ODE Formulation}

\begin{theorem}[Continuous Limit of the Difference Equation]
The first-order condition $\nabla \tilde{U}(x^*) = 0 \implies A x^* = b$ implies that for the interior points $1 < k < N-1$, the sequence satisfies:
\[ \frac{x_{k+1} - 2x_k + x_{k-1}}{\tau^2} = \frac{\lambda \sigma^2}{\tilde{\eta}} x_k \]
Taking the limit as $\tau \to 0$, this difference equation converges to the ordinary differential equation:
\[ x''(t) - \tilde{\kappa}^2 x(t) = 0 \]
where $\tilde{\kappa} = \sigma \sqrt{\frac{\lambda}{\tilde{\eta}}}$.
\end{theorem}

\begin{proof}
Extracting the $k$-th row from the linear system $Ax = b$ for an interior point gives:
\[ -\frac{2\tilde{\eta}}{\tau} x_{k-1} + \left(\frac{4\tilde{\eta}}{\tau} + 2\lambda\sigma^2\tau\right) x_k - \frac{2\tilde{\eta}}{\tau} x_{k+1} = 0 \]
Dividing by $-\frac{2\tilde{\eta}}{\tau}$ and rearranging the terms yields:
\[ x_{k+1} - 2x_k + x_{k-1} = \frac{\lambda \sigma^2 \tau^2}{\tilde{\eta}} x_k \]
Dividing both sides by $\tau^2$ produces the discrete second derivative operator on the left side. As $\tau \to 0$, the left side formally converges to $x''(t)$, yielding $x''(t) - \tilde{\kappa}^2 x(t) = 0$.
\end{proof}

\begin{proposition}[Exact Analytical Solution]
The continuous-time ODE subject to the boundary conditions $x(0) = X$ and $x(T) = 0$ admits the closed-form analytical solution:
\[ x(t) = X \frac{\sinh(\tilde{\kappa}(T-t))}{\sinh(\tilde{\kappa} T)} \]
\end{proposition}

\begin{proof}
The characteristic equation for the ODE is $r^2 - \tilde{\kappa}^2 = 0$, yielding roots $r = \pm \tilde{\kappa}$. The general solution in exponential basis is:
\[ x(t) = A e^{\tilde{\kappa}t} + B e^{-\tilde{\kappa}t} \]
Applying the boundary conditions $x(0) = X$ and $x(T) = 0$ yields a system of equations for the constants, resolving to:
\[ A = \frac{-e^{-\tilde{\kappa}T}}{2\sinh(\tilde{\kappa}T)}X, \quad B = \frac{e^{\tilde{\kappa}T}}{2\sinh(\tilde{\kappa}T)}X \]
Substituting these constants back into the general solution:
\[ x(t) = \frac{X}{2\sinh(\tilde{\kappa}T)} \left( e^{\tilde{\kappa}(T-t)} - e^{-\tilde{\kappa}(T-t)} \right) \]
By the definition of the hyperbolic sine, $\sinh(z) = \frac{e^z - e^{-z}}{2}$, this simplifies strictly to:
\[ x(t) = X \frac{\sinh(\tilde{\kappa}(T-t))}{\sinh(\tilde{\kappa}T)} \]
\end{proof}

\begin{remark}
The continuous model confirms that the discrete matrix mechanics align strictly with the market dynamics. The structural constant $\tilde{\kappa}$ governs the curvature of the liquidation scheme; a higher $\tilde{\kappa}$ implies higher execution urgency. Numerically, evaluating the continuous solution against the discrete $O(N)$ solver yields deviations on the order of $0.11\%$, isolated entirely to discretization error.
\end{remark}

\section{Empirical Validation and Efficient Frontier Mechanics}

To rigorously validate the discrete mathematical filtrations of the objective function, an event-driven Monte Carlo engine was constructed to simulate $M=10,000$ price paths. 

\paragraph{Implied Parameter Calibration}
Due to the absence of proprietary limit order book data, the market impact parameters were determined via implied calibration. The parameters were algebraically reverse-engineered to ensure the risk-neutral TWAP strategy yields a standard institutional friction cost of roughly 25 basis points (bps) for a $\$1,000,000$ mid-cap equity position ($X = 10,000$ shares, $S_0 = 100$). 

Setting the temporary impact $\eta = 2 \times 10^{-5}$ and permanent impact $\gamma = 1 \times 10^{-5}$ over $N=50$ intervals ($\tau=0.02$) anchors the baseline TWAP execution cost to exactly $\$2,490$ (24.9 bps). The arithmetic volatility is set to $\sigma = 0.3$. This scales the maximum time-exposure risk to approximately 17 bps, precisely reflecting the realistic standard deviation of an intraday execution.

\begin{proposition}[Maximum Variance Limit]
The physical boundary of the Efficient Frontier is strictly bounded on the right by the risk-neutral TWAP strategy ($\lambda \to 0$). The absolute maximum variance exposure is given by:
\[ V_{max} = \sigma^2 \tau \sum_{j=1}^{N-1} \left( X \frac{j}{N} \right)^2 \]
\end{proposition}

\begin{proof}
By substituting the deterministic TWAP inventory schedule $x_k = X \frac{N-k}{N}$ into the variance equation $V[IS] = \sigma^2 \tau \sum_{k=1}^{N-1} x_k^2$, we obtain:
\[ V_{max} = \sigma^2 \tau \sum_{k=1}^{N-1} \left( X \frac{N-k}{N} \right)^2 \]
Changing the summation index to $j = N-k$ reverses the sum but preserves the terms precisely, yielding the result. The variance cannot physically exceed this value, as TWAP represents the maximum possible time-exposure of the inventory to stochastic price shocks. With our calibrated parameters, $V_{max} \approx 2.91 \times 10^6$, corresponding to an execution standard deviation of $\$1,706$ ($17$ bps).
\end{proof}

\paragraph{Structural Bounds vs. The Institutional Operating Window}
To understand the global properties of the continuous ODE versus the local constraints of a trading desk, we isolate the Efficient Frontier into two distinct scopes. 

Figure \ref{fig:theoretical_bounds} displays the macro-mathematical view ($\lambda \in [10^{-6}, 10^1]$). As $\lambda \to \infty$, the variance penalty forces a corner solution: an instantaneous block liquidation ($v_1 \to X$). Under the infinite liquidity assumption of the linear framework, this yields a theoretical shortfall of approximately $-\$560,000$. In practice, this 56\% slippage represents a physical liquidity wall---a mathematically valid limit that no rational portfolio manager would ever execute.

\begin{figure}[H]
    \centering
    \includegraphics[width=0.8\textwidth]{results/Efficient_Frontier_high_range.png}
    \caption{Theoretical Bounds ($\lambda \in [10^{-6}, 10^1]$). The left boundary represents the prohibitive market impact of an instantaneous block liquidation.}
    \label{fig:theoretical_bounds} 
\end{figure}

Figure \ref{fig:practical_bounds} restricts the domain to the viable operational window ($\lambda \in [10^{-6}, 10^{-1}]$). Here, execution costs range from the 24.9 bps TWAP baseline up to approximately 200 bps. Quantitative funds restrict their optimization entirely to this flatter sub-domain. The exact chosen coordinate balances minor temporary impact penalties against strategy-specific constraints, such as alpha decay (directional drift), preserving tight Sharpe ratios in statistical arbitrage, or preempting volatility shocks.

\begin{figure}[H]
    \centering
    \includegraphics[width=0.8\textwidth]{results/Efficient_Frontier_low_range.png}
    \caption{The Rational Execution Window ($\lambda \in [10^{-6}, 10^{-1}]$). The strict operational domain where execution algorithms balance impact costs against standard intraday variance.}
    \label{fig:practical_bounds} 
\end{figure}

\paragraph{Invariance to Permanent Impact}
While the permanent impact parameter $\gamma$ explicitly penalizes the total expected shortfall (shifting the Efficient Frontier vertically along the y-axis), it entirely vanishes from the Hessian matrix $A$ and the linear vector $b$ in the quadratic form $\tilde{U}(x)$. Mathematically, $\gamma$ plays absolutely no role in shaping the optimal execution trajectory $x^*$.

\section{Machine Learning Engine and Predictive Microstructure}

To extend the static model into a dynamic Model Predictive Control (MPC) framework, we integrate localized order book features (Garman-Klass volatility, Order Flow Imbalance, Market Depth) to predict short-term regime shifts in volatility ($\sigma$) and liquidity ($\eta$).

\paragraph{The Positive-Definite Matrix Constraint}
The discrete Almgren-Chriss solver requires the symmetric tridiagonal matrix $A$ to be strictly positive-definite to guarantee a unique global minimum. This requires the condition:
\[ \tilde{\eta} = \eta - \frac{1}{2}\gamma\tau > 0 \]
Standard linear regressors can output negative predictions, which mathematically breaks the solver. To enforce this strict physical boundary, the target variables are transformed into log-space ($\ln(\sigma)$, $\ln(\eta)$) prior to model training. The resulting predictions are exponentiated ($e^{\text{pred}}$) back to real-space before being passed to the dynamic execution engine, guaranteeing strictly positive parameters.

\paragraph{Linear Limitations and the Non-Linear Pivot}
Initial out-of-sample testing utilized a Ridge regressor with dynamic $\alpha$ tuning. When benchmarked against a naive 10-period Simple Moving Average (SMA), the Ridge model outperformed on liquidity prediction (MSE improvement of $+4.90\%$) but suffered a massive degradation in volatility prediction ($-31.04\%$ against the SMA). 

This empirical failure reveals a core microstructural reality: order book dynamics are highly non-linear. Market phenomena such as sudden liquidity vacuums or volatility spikes are threshold events. Continuous linear weights fail to capture these regime shifts, capturing noise instead of signal. Consequently, the architecture was pivoted to tree-based non-linear regressors (Random Forest), which inherently map step-function thresholds (e.g., if depth drops below $x$ and momentum exceeds $y$, apply a discrete volatility spike).

\section{Empirical Bootstrapping and Dynamic MPC Validation}

Theoretical validation of the Almgren-Chriss framework assumes a zero-drift Martingale, which mathematically guarantees a strictly negative Implementation Shortfall (IS). However, to rigorously evaluate the dynamic MPC logic under real-world conditions, we constructed an empirical bootstrapper to sample $M=1,500$ historical execution windows of AAPL.

\paragraph{The Evaluation Engine Framework}
Before comparing the static and dynamic trajectories, the historical grading engine required two fundamental structural adjustments to ensure mathematical fairness:
\begin{enumerate}
    \item \textbf{Cumulative Permanent Impact:} Raw historical prices represent a universe where the trades did not occur. To evaluate the strategy accurately, the cumulative permanent impact penalty ($\gamma$) was iteratively subtracted from the historical price path to simulate the actual market depression caused by our liquidations.
    \item \textbf{Localized Liquidity Scoring:} Evaluating the temporary impact using a flat historical average ($\bar{\eta}$) mathematically blinds the evaluator to the MPC's primary objective. The bootstrapper was updated to score execution costs against the localized $\eta_t$ arrays. This ensures the Machine Learning model is properly rewarded for routing volume away from predicted liquidity droughts.
\end{enumerate}

\paragraph{The Noise-to-Impact Ratio: A Mathematical Breakdown}
When executing a highly liquid mega-cap equity, the theoretical execution penalty is frequently overpowered by the asset's intraday Brownian variance. We can decompose the exact orders of magnitude for our historical AAPL simulation:

\begin{itemize}
    \item \textbf{Initial Position Value ($X S_0$):} Liquidating $X = 50,000$ shares at an arrival price of $S_0 \approx \$335$ yields a total portfolio exposure of roughly $\$16,750,000$.
    \item \textbf{Market Variance ($\sigma$):} The standard deviation of the position value over the intraday execution horizon is approximately $\$35,000$. This represents the absolute market noise.
    \item \textbf{Execution Cost ($\eta, \gamma$):} For a risk-neutral TWAP schedule, the volume per step is $v_k = 1,000$ shares with a time step of $\tau = 0.02$. Using the historically realized temporary impact of $\eta \approx 1.7 \times 10^{-9}$, the total temporary impact cost evaluates to strictly pennies: $\frac{\eta}{\tau} \sum v_k^2 \approx \$4.25$.
\end{itemize}

Because the fundamental market impact ($\$4.25$) is microscopically small compared to the intraday market noise ($\pm \$35,000$), the final IS is entirely dictated by directional drift. During the sampled historical period, AAPL exhibited a positive trend. This natural market drift generated more cash than the arrival price benchmark, shifting the expected IS mathematically into the positive domain ($E[IS] \approx +\$2,445$).

\begin{figure}[H]
    \centering
    \includegraphics[width=0.8\textwidth]{results/Empirical_IS_Distribution_AAPL_Base.png}
    \caption{Empirical IS Distribution ($M=1,500$). The expected shortfall rests in the positive domain ($+\$2,445$) because the natural historical market drift outpaced the negligible $\$4.25$ market impact of a highly liquid asset.}
    \label{fig:empirical_is_base} 
\end{figure}

\paragraph{Mathematical Validation via Illiquidity Scaling}
To prove that the MPC grading engine and the discrete tridiagonal matrix solver are mathematically sound despite the positive baseline IS, we must isolate the Almgren-Chriss penalty from the market noise. 

We conducted a stress test by artificially scaling the impact parameters ($\eta$ and $\gamma$) by a factor of $10^4$. This simulation transforms the highly liquid AAPL data into a severely illiquid micro-cap environment, inflating the baseline execution costs to approximately $\$400,000$, thereby aggressively overpowering the $\$35,000$ intraday variance. 

\begin{figure}[H]
    \centering
    \includegraphics[width=0.8\textwidth]{results/Empirical_IS_Distribution_Scaled_1e4.png}
    \caption{Stress-Tested IS Distribution with impact parameters scaled by $10^4$. Stripping away the Brownian noise forces the empirical distribution strictly below zero, mathematically validating the discrete execution penalty.}
    \label{fig:empirical_is_scaled} 
\end{figure}

As demonstrated in Figure \ref{fig:empirical_is_scaled}, artificially suppressing the noise-to-impact ratio forces the entire distribution strictly into the negative domain. The expected Implementation Shortfall collapses to approximately $-\$2,000,000$. The near-perfect alignment between the static and dynamic expected shortfalls in this regime confirms a critical theoretical boundary: when liquidity constraints become infinitely severe, the dynamic MPC solver correctly suppresses all opportunistic machine learning signals and flattens out, defaulting strictly back to the optimal variance-minimizing schedule.

\end{document}